\documentclass[12pt, a4paper]{article}

% ── PACKAGES ──────────────────────────────────────────────────
\usepackage[utf8]{inputenc}
\usepackage[T1]{fontenc}
\usepackage{geometry}
\geometry{margin=2.5cm}
\usepackage{hyperref}
\hypersetup{
    colorlinks=true,
    linkcolor=blue,
    urlcolor=blue,
    citecolor=blue,
    pdftitle={A Six Lock Type Classification of Breast
              Cancer: Waddington Attractor Geometry Applied
              to All Major Subtypes},
    pdfauthor={Eric Robert Lawson},
    pdfkeywords={breast cancer, attractor geometry,
                 Waddington landscape, lock type,
                 epigenetic classification, FOXA1, EZH2,
                 TNBC, LumA, LumB, HER2, claudin-low,
                 ILC, CDK4/6, HDACi, tazemetostat,
                 OrganismCore, cancer taxonomy}
}
\usepackage{amsmath}
\usepackage{booktabs}
\usepackage{longtable}
\usepackage{array}
\usepackage{parskip}
\usepackage{microtype}
\usepackage{authblk}
\usepackage{titlesec}
\usepackage{fancyhdr}
\usepackage{xcolor}
\usepackage{float}
\usepackage{enumitem}
\usepackage{setspace}
\usepackage{multirow}

% ── COLOURS ───────────────────────────────────────────────────
\definecolor{repobox}{RGB}{240,247,255}
\definecolor{lockbox}{RGB}{245,255,245}
\definecolor{mechanismbox}{RGB}{250,245,255}
\definecolor{novelbox}{RGB}{255,248,220}
\definecolor{summarybox}{RGB}{240,255,240}
\definecolor{luma}{RGB}{65,105,225}
\definecolor{lumb}{RGB}{30,144,255}
\definecolor{her2}{RGB}{255,140,0}
\definecolor{tnbc}{RGB}{220,20,60}
\definecolor{cl}{RGB}{128,0,128}
\definecolor{ilc}{RGB}{34,139,34}

% ── HEADER / FOOTER ───────────────────────────────────────────
\pagestyle{fancy}
\fancyhf{}
\rhead{\small OrganismCore \textbar{} 2026-03-06}
\lhead{\small CS-LIT-2: The Six Lock Type Classification}
\rfoot{\small Page \thepage}
\renewcommand{\headrulewidth}{0.4pt}

% ── SECTION FORMATTING ─────────────────────────────────��──────
\titleformat{\section}
  {\large\bfseries}{\thesection}{1em}{}[\titlerule]
\titleformat{\subsection}
  {\normalsize\bfseries}{\thesubsection}{1em}{}
\titleformat{\subsubsection}
  {\small\bfseries}{\thesubsubsection}{1em}{}

% ══════════════════════════════════════════════════════════════
% TITLE
% ══════════════════════════════════════════════════════════════
\title{
    \vspace{-1em}
    \textbf{A Six Lock Type Classification of Breast
    Cancer:}\\[0.5em]
    \large Waddington Attractor Geometry Identifies
    Distinct Epigenetic and Structural Locks in All
    Major Subtypes, Each with a Specific Therapeutic
    Unlock Sequence\\[0.4em]
    \normalsize \textit{OrganismCore Framework ---
    Document CS-LIT-2}\\[0.3em]
    \small Derived from First Principles Before
    Literature Review\\[0.3em]
    \small Framework Applied to 19,542 Single Cancer
    Cells Across Six Subtypes (GSE176078)
}

\author[1]{Eric Robert Lawson}
\affil[1]{
    OrganismCore\\
    \href{mailto:OrganismCore@proton.me}
         {OrganismCore@proton.me}\\
    ORCID:
    \href{https://orcid.org/0009-0002-0414-6544}
         {0009-0002-0414-6544}
}

\date{March 6, 2026}

% ══════════════════════════════════════════════════════════════
\begin{document}
% ══════════════════════════════════════════════════════════════

\maketitle
\thispagestyle{fancy}

% ── REPOSITORY POINTER ────────────────────────────────────────
\begin{center}
\colorbox{repobox}{
\begin{minipage}{0.88\textwidth}
\vspace{0.4em}
\centering
\textbf{Full computational record, scripts, reasoning
artifacts, and subtype deep dives:}\\[0.3em]
\href{https://github.com/Eric-Robert-Lawson/attractor-oncology}
{\texttt{https://github.com/Eric-Robert-Lawson/attractor-oncology}}
\\[0.2em]
\small
Repository ID: 1172048282 \textbar{}
MIT License \textbar{}
Python \textbar{}
Public\\[0.2em]
Source document: \texttt{Cancer\_Research/BRCA/DEEP\_DIVE/}
\texttt{BRCA\_Cross\_Subtype\_Literature\_Check.md}
(BRCA-S8h, item CS-LIT-2)\\[0.2em]
Subtype deep dives: \texttt{LumA/}, \texttt{LumB/},
\texttt{HER2\_ENRICHED/}, \texttt{TNBC/},
\texttt{Claudin-low/}, \texttt{ILC/}\\[0.2em]
Companion deposits: CS-LIT-1 (FOXA1/EZH2 ratio),
CS-LIT-16 (tazemetostat $\rightarrow$ fulvestrant),
CS-LIT-17 (tazemetostat maintenance)
--- all Zenodo 2026-03-06
\vspace{0.4em}
\end{minipage}
}
\end{center}

\vspace{0.8em}

% ── ABSTRACT ──────────────────────────────────────────────────
\begin{abstract}
\noindent
Breast cancer is routinely described as a collection
of molecularly distinct subtypes. What is not yet
described is why those subtypes exist as a coherent
ordered set, what single biological axis organises
them, and --- critically --- what specific molecular
lock maintains each subtype in its aberrant state
and what specific therapeutic action dissolves that
lock. This document presents the OrganismCore
six lock type classification of breast cancer,
derived from Waddington landscape attractor geometry
applied to 19,542 single cancer cells from 26
patients across six subtypes (GSE176078, Wu et al.\
\textit{Nature Genetics} 2021) before any literature
review was performed.

The classification identifies six distinct breast
cancer lock types, each corresponding to a specific
attractor geometry, a specific molecular maintenance
mechanism, and a specific therapeutic unlock sequence:
(1)~\textbf{Luminal A} --- CDK4/6-mediated cycle lock
with intact luminal identity (Type~1, slope arrest);
(2)~\textbf{Luminal B} --- DNMT3A/HDAC2 chromatin
lock suppressing ER output (Type~1, slope arrest with
epigenetic brake); (3)~\textbf{HER2-enriched} ---
ERBB2 amplicon signalling override of the luminal
slope (Type~1, kinase arrest);
(4)~\textbf{TNBC/Basal-like} --- EZH2/PRC2
H3K27me3 epigenetic lock silencing luminal identity
(composite Type~1$\rightarrow$2, wrong valley);
(5)~\textbf{Claudin-low} --- stem cell root lock
with no residual luminal programme
(Type~4, root lock); (6)~\textbf{Invasive Lobular
Carcinoma (ILC)} --- CDH1 structural lock creating
the geometric inverse of TNBC (Type~3, structural
lock).

Each lock type predicts a specific first therapeutic
action before any subsequent agent can engage. The
six lock types are ordered by a single continuous
variable --- the FOXA1/EZH2 ratio --- which encodes
the balance between luminal identity maintenance and
epigenetic identity suppression. This ordering was
confirmed across $\sim$7,500 patients in seven
independent datasets on four measurement platforms
(companion Zenodo deposit, 2026-03-06). The
classification provides a mechanistically unified
explanation for why these six subtypes exist, why
they respond differently to the same agents, and why
each requires a specific first step that is unique
to its lock type. No competing classification
framework using attractor geometry as its organising
principle currently exists in the published literature.
\end{abstract}

\vspace{0.5em}
\noindent\textbf{Keywords:}
breast cancer; attractor geometry; Waddington landscape;
lock type; epigenetic classification; FOXA1; EZH2;
luminal A; luminal B; HER2-enriched; TNBC; basal-like;
claudin-low; invasive lobular carcinoma; CDK4/6;
HDACi; entinostat; tazemetostat; fulvestrant; CDH1;
PRC2; H3K27me3; DNMT3A; HDAC2; OrganismCore;
cancer taxonomy; precision oncology; drug sequencing

\newpage
\tableofcontents
\newpage

% ══════════════════════════════════════════════════════════════
\section{Document Metadata}
% ══════════════════════════════════════════════════════════════

\begin{center}
\begin{tabular}{ll}
\toprule
\textbf{Field} & \textbf{Value} \\
\midrule
Document ID     & CS-LIT-2-SIX-LOCK-CLASSIFICATION \\
Series          & BRCA Deep Dive --- Cross-Subtype \\
Type            & Framework Classification Paper ---
                  Priority Establishment Document \\
Parent document & BRCA-S8h (Cross-Subtype Literature
                  Check, item CS-LIT-2) \\
Derivation data & GSE176078 (Wu et al.\ 2021,
                  \textit{Nature Genetics}) \\
Validation data & TCGA-BRCA, METABRIC, CPTAC-BRCA,
                  GSE96058, GSE25066 ($\sim$7,500 patients) \\
Companion deposits
& CS-LIT-1 FOXA1/EZH2 ratio (Zenodo 2026-03-06)\\
& CS-LIT-16 tazemetostat $\rightarrow$ fulvestrant\\
& CS-LIT-17 tazemetostat maintenance \\
Date            & 2026-03-06 \\
Author          & Eric Robert Lawson \\
Organisation    & OrganismCore \\
ORCID           & 0009-0002-0414-6544 \\
Repository      &
  \href{https://github.com/Eric-Robert-Lawson/attractor-oncology}
       {Eric-Robert-Lawson/attractor-oncology} \\
Repository ID   & 1172048282 \\
License         & CC BY 4.0 \\
Biological contradictions & 0 \\
\bottomrule
\end{tabular}
\end{center}

% ══════════════════════════════════════════════════════════════
\section{The Organising Principle: Waddington Landscape
Attractor Geometry}
% ══════════════════════════════════════════════════════════════

\subsection{The Waddington model applied to cancer}

The Waddington developmental landscape represents cell
differentiation as a ball rolling down a terrain of
valleys and ridges. Each valley (attractor) corresponds
to a stable cell identity state. Committed cells roll
into their correct terminal attractor and remain there.
Stem and progenitor cells are at the top of the
landscape --- they have not yet committed and can roll
into multiple possible valleys.

Cancer arises when cells are arrested in abnormal
attractor states. The OrganismCore framework identifies
four canonical attractor geometry types:

\begin{center}
\begin{tabular}{p{1.5cm}p{3cm}p{9cm}}
\toprule
\textbf{Type} & \textbf{Name} & \textbf{Geometry} \\
\midrule
Type 1
& Slope Arrest
& Cell arrested on the slope approaching the correct
  terminal attractor. It has not reached its normal
  identity but has not fallen into a wrong valley.
  The correct identity programme is partially active.
  A molecular brake is preventing completion. \\[0.3em]
Type 2
& Wrong Valley
& Cell has fallen into the wrong attractor. It is
  expressing the identity of a different cell type.
  The correct identity programme is silenced.
  An epigenetic lock maintains the wrong state. \\[0.3em]
Type 3
& Structural Lock
& Cell is held in an arrested state by a structural
  (non-epigenetic) modification --- loss of a cell
  adhesion or cytoskeletal component that normally
  constrains cell identity. \\[0.3em]
Type 4
& Root Lock
& Cell is arrested at the progenitor root. It has
  never committed to any terminal identity.
  The full differentiation programme is absent,
  not silenced. \\
\bottomrule
\end{tabular}
\end{center}

\subsection{The single axis that orders all six subtypes}

All six breast cancer subtypes are ordered by a single
continuous variable: the FOXA1/EZH2 ratio. This variable
measures the balance between luminal identity presence
(FOXA1, pioneer factor that opens and maintains luminal
chromatin) and luminal identity suppression (EZH2,
Polycomb repressor that silences the luminal programme
by depositing H3K27me3).

\begin{center}
\begin{tabular}{lccc}
\toprule
\textbf{Subtype} & \textbf{FOXA1/EZH2 ratio}
& \textbf{Lock type} & \textbf{Attractor type} \\
\midrule
Luminal A     & 9.38 & CDK4/6 cycle lock & Type 1 \\
Luminal B     & 8.10 & DNMT3A/HDAC2 chromatin lock
                     & Type 1 \\
HER2-enriched & 3.34 & ERBB2 kinase override
                     & Type 1 \\
TNBC/Basal    & 0.52 & EZH2/PRC2 epigenetic lock
                     & Composite 1$\rightarrow$2 \\
Claudin-low   & 0.10 & Stem root lock (no programme)
                     & Type 4 \\
ILC           & ---  & CDH1 structural lock
                     & Type 3 \\
\bottomrule
\end{tabular}
\end{center}

\noindent ILC is not placed on the FOXA1/EZH2 ratio axis
because its lock is structural, not epigenetic. It is the
geometric inverse of TNBC: where TNBC has an epigenetic
lock (EZH2 silences the luminal programme), ILC has a
structural lock (CDH1 loss removes the adhesion
constraint that keeps luminal cells organised). Both
are derangements of luminal identity, but from opposite
directions.

% ══════════════════════════════════════════════════════════════
\section{Lock Type 1: Luminal A --- The CDK4/6 Cycle Lock}
% ══════════════════════════════════════════════════════════════

\subsection{Attractor geometry}

\begin{center}
\colorbox{mechanismbox}{
\begin{minipage}{0.85\textwidth}
\vspace{0.5em}
\textbf{LUMINAL A --- TYPE 1 SLOPE ARREST}

\medskip
\textbf{Cell of origin:} Mature luminal epithelial cell\\
\textbf{Normal terminal attractor:} ER-positive,
PR-positive, low-proliferation luminal cell\\
\textbf{Arrest point:} On the luminal slope, with full
luminal identity intact but cell cycle dysregulated\\
\textbf{Lock:} CDK4/6-mediated cycle entry following
CDKN1A (p21) reduction\\
\textbf{FOXA1/EZH2 ratio:} 9.38 (highest of all subtypes)\\
\textbf{Luminal programme status:} Fully intact

\medskip
\textbf{Key geometry:}\\
FOXA1 $+$highest, ESR1 $+$highest, GATA3 $+$highest\\
EZH2 $-$lowest\\
The cell is fully luminal. The problem is not identity
suppression. The problem is cell cycle escape.
\vspace{0.5em}
\end{minipage}
}
\end{center}

\subsection{The molecular lock}

LumA cells have:
\begin{itemize}[noitemsep]
  \item Intact luminal identity programme (FOXA1, GATA3,
        ESR1, PGR all high)
  \item Reduced CDKN1A (p21) expression --- the
        CDK4/6 brake is weakened
  \item RB1 hyperphosphorylation by CDK4/6 ---
        cell cycle entry proceeds despite low proliferation
        rate in absolute terms
  \item Low EZH2 (lowest of all subtypes) ---
        the silencing programme is not engaged
\end{itemize}

The lock is in the cell cycle compartment, not in the
identity compartment. The cell knows what it is.
It is cycling when it should not be.

\subsection{The therapeutic unlock}

\begin{center}
\begin{tabular}{p{3.5cm}p{9cm}}
\toprule
\textbf{Step} & \textbf{Action and rationale} \\
\midrule
First line
& \textbf{CDK4/6 inhibitor} (palbociclib, ribociclib,
  or abemaciclib) $+$ endocrine therapy (aromatase
  inhibitor or fulvestrant). The CDK4/6 inhibitor
  directly addresses the lock: restores RB1
  hypophosphorylation, re-establishes cell cycle arrest.
  Endocrine therapy engages the intact ER simultaneously. \\[0.3em]
Biomarker (novel)
& \textbf{CDKN1A (p21) level as quantitative
  predictor of CDK4/6i benefit magnitude.} Lower
  p21 at baseline predicts greater magnitude of
  CDK4/6i benefit because the CDK4/6 brake was more
  severely weakened. This is a continuous predictor,
  not a binary classifier. Testable in PALOMA archived
  tissue without new trial. (CS-LIT-14) \\
\bottomrule
\end{tabular}
\end{center}

\subsection{Published confirmation}

CDK4/6 inhibitors are standard of care in LumA/HR+
breast cancer. The geometry prediction (cycle lock in
intact luminal cell) is fully convergent with the
clinical development of palbociclib, ribociclib, and
abemaciclib. The novel contribution is the CDKN1A
quantitative predictor (CS-LIT-14), which is not in
the published literature as a continuous variable.

% ══════════════════════════════════════════════════════════════
\section{Lock Type 2: Luminal B --- The DNMT3A/HDAC2
Chromatin Lock}
% ══════════════════════════════════════════════════════════════

\subsection{Attractor geometry}

\begin{center}
\colorbox{mechanismbox}{
\begin{minipage}{0.85\textwidth}
\vspace{0.5em}
\textbf{LUMINAL B --- TYPE 1 SLOPE ARREST WITH
EPIGENETIC BRAKE}

\medskip
\textbf{Cell of origin:} Mature luminal epithelial cell\\
\textbf{Normal terminal attractor:} ER-positive,
low-proliferation luminal cell\\
\textbf{Arrest point:} On the luminal slope with
luminal identity present but ER output suppressed
by chromatin remodelling\\
\textbf{Lock:} DNMT3A/HDAC2 co-expression coupling
suppresses TFF1/ESR1 transcriptional output\\
\textbf{FOXA1/EZH2 ratio:} 8.10 (high, just below LumA)\\
\textbf{Luminal programme status:} Present but muted

\medskip
\textbf{Key geometry:}\\
FOXA1 high (but lower than LumA)\\
DNMT3A/HDAC2 co-expression coupling: r$=+0.267$ in LumB
vs r$=+0.071$ in LumA ($p=5.68\times10^{-56}$)\\
TFF1/ESR1 decoupled: r lower in LumB than LumA
($p=0.0019$ in METABRIC replication)\\
The cell has luminal identity. The ER circuit is
present but the chromatin state is suppressing
its output.
\vspace{0.5em}
\end{minipage}
}
\end{center}

\subsection{The molecular lock}

LumB is distinguished from LumA by a specific
epigenetic feature: DNMT3A and HDAC2 are
co-expressed at levels that produce active
chromatin compaction at ER-target gene promoters.
The consequences:

\begin{itemize}[noitemsep]
  \item TFF1 (an ER-target gene and ER circuit
        output marker) is decoupled from ESR1 ---
        the ER is present but its transcriptional
        output is suppressed
  \item This suppression is HDAC2-mediated chromatin
        compaction, not EZH2-mediated H3K27me3
        (EZH2 is slightly elevated in LumB but is
        not the primary lock)
  \item The DNMT3A/HDAC2 co-expression coupling
        ($r=+0.267$ in LumB vs $r=+0.071$ in LumA)
        is the specific geometric signature of
        the LumB lock
  \item Endocrine therapy alone is insufficient
        because the ER is present but its chromatin
        is not fully accessible to downstream
        transcriptional targets
\end{itemize}

\subsection{The therapeutic unlock}

\begin{center}
\begin{tabular}{p{3.5cm}p{9cm}}
\toprule
\textbf{Step} & \textbf{Action and rationale} \\
\midrule
First step
& \textbf{HDAC inhibitor} (entinostat specifically ---
  class I HDAC, HDAC1/2/3 selective) to open the
  compacted chromatin at ER-target gene promoters.
  This directly addresses the HDAC2-mediated lock. \\[0.3em]
Second step
& \textbf{Fulvestrant} (or aromatase inhibitor) to
  engage the ER once its chromatin context is restored.
  Without the HDACi first, endocrine therapy cannot
  fully engage because the downstream targets are
  chromatinised shut. \\[0.3em]
Patient selector (novel)
& \textbf{TFF1/ESR1 ratio as LumB patient selector}
  for HDACi benefit. Low TFF1/ESR1 (decoupled)
  = chromatin lock present = HDACi benefit predicted.
  High TFF1/ESR1 (coupled) = no lock = no additional
  HDACi benefit. Testable in NCT07235618 (actively
  recruiting, January 2026). (CS-LIT-23) \\
\bottomrule
\end{tabular}
\end{center}

\subsection{Published confirmation}

Entinostat $+$ exemestane: ENCORE 301 showed OS
benefit ($p=0.036$) in ER$+$ patients. The E2112
trial showed no benefit in unselected ER$+$ patients.
The framework's contribution is the resolution of
this contradiction: entinostat benefit is specifically
in the DNMT3A/HDAC2-locked LumB subgroup, not in
all ER$+$ patients. The TFF1/ESR1 ratio (CS-LIT-23)
is the patient selector that identifies this subgroup.

% ══════════════════════════════════════════════════════════════
\section{Lock Type 3: HER2-Enriched --- The ERBB2
Kinase Override}
% ══════���═══════════════════════════════════════════════════════

\subsection{Attractor geometry}

\begin{center}
\colorbox{mechanismbox}{
\begin{minipage}{0.85\textwidth}
\vspace{0.5em}
\textbf{HER2-ENRICHED --- TYPE 1 SLOPE ARREST,
KINASE OVERRIDE}

\medskip
\textbf{Cell of origin:} Luminal progenitor
(luminal epithelial cell)\\
\textbf{Normal terminal attractor:} Mature luminal cell\\
\textbf{Arrest point:} On the luminal slope, arrested
by ERBB2 amplicon signalling override\\
\textbf{Lock:} ERBB2 amplification on chromosome 17q12
drives constitutive PI3K/MAPK signalling, overriding
normal luminal differentiation signals\\
\textbf{FOXA1/EZH2 ratio:} 3.34 (intermediate)\\
\textbf{Luminal programme status:} Partially retained
(FOXA1 elevated $+12.7\%$, GATA3 present $-44.7\%$,
ESR1 near-absent $-92.5\%$, PGR near-absent $-95.2\%$)

\medskip
\textbf{Key geometry:}\\
The HER2 cell is not in the wrong valley (TYPE 2 excluded
--- no SOX10, no KRT5/14, no VIM).\\
It is on the luminal slope but the signalling that would
complete differentiation is overwhelmed by the amplicon.\\
This is TYPE 1 with a signalling-level brake, not an
epigenetic brake.
\vspace{0.5em}
\end{minipage}
}
\end{center}

\subsection{The molecular lock}

The 17q12 amplicon encoding ERBB2 (and co-amplified
STARD3, GRB7, MIEN1) drives:
\begin{itemize}[noitemsep]
  \item Constitutive PI3K/AKT/mTOR activation
  \item Constitutive MAPK/ERK activation
  \item Suppression of ESR1 and PGR transcription
        (ER circuit severed despite retained FOXA1)
  \item High proliferation (Ki67 elevated)
  \item Resistance to endocrine therapy alone
        (ER circuit severed by signalling override)
\end{itemize}

The lock is not epigenetic in origin. The chromatin
is not the primary problem. The kinase signalling
cascade is the lock. This is why neither CDK4/6
inhibitors (cycle lock absent), HDACi (chromatin
lock absent), nor EZH2 inhibitors (EZH2 is moderately
elevated but is not the primary lock) are the correct
first step. Anti-HER2 is the correct first step
because it addresses the lock directly.

\subsection{The therapeutic unlock}

\begin{center}
\begin{tabular}{p{3.5cm}p{9cm}}
\toprule
\textbf{Step} & \textbf{Action and rationale} \\
\midrule
First step
& \textbf{Anti-HER2 therapy} (trastuzumab,
  pertuzumab, T-DM1, or tucatinib depending on line)
  to remove the amplicon signalling override. \\[0.3em]
Second step
& \textbf{Endocrine therapy} for the ER$+$/HER2$+$
  (luminal/HER2) fraction once the amplicon override
  is suppressed. The retained FOXA1 allows partial
  ER circuit re-engagement after HER2 suppression. \\[0.3em]
Subgroup (novel)
& \textbf{EZH2i $+$ anti-HER2 for the HER2-deep
  fraction:} CDH3-high, AR-low, EZH2$+118\%$
  HER2-enriched patients have a secondary epigenetic
  component. This subgroup requires EZH2 inhibition
  in addition to anti-HER2. (CS-LIT-21) \\
\bottomrule
\end{tabular}
\end{center}

\subsection{Published confirmation}

Anti-HER2 therapy is established standard of care.
The geometry prediction (kinase arrest on luminal
slope, not wrong valley) is confirmed by the
observation that HER2-enriched tumours retain FOXA1
and respond to endocrine therapy after HER2
suppression. The HER2-deep subgroup (CS-LIT-21)
with secondary EZH2 lock is the novel contribution.

% ══════════════════════════════════════════════════════════════
\section{Lock Type 4: TNBC/Basal-like --- The EZH2/PRC2
Epigenetic Lock}
% ══════════════════════════════════════════════════════════════

\subsection{Attractor geometry}

\begin{center}
\colorbox{mechanismbox}{
\begin{minipage}{0.85\textwidth}
\vspace{0.5em}
\textbf{TNBC/BASAL-LIKE --- COMPOSITE TYPE 1
$\rightarrow$ TYPE 2}

\medskip
\textbf{Cell of origin:} BRCA1-regulated luminal
progenitor\\
\textbf{Normal terminal attractor:} Mature luminal cell\\
\textbf{Arrest mechanism:} Two-stage composite:\\
\quad Stage 1 (Type 1): BRCA1 loss blocks luminal
progenitor differentiation (blocked approach)\\
\quad Stage 2 (Type 2): Progenitor falls into
basal/myoepithelial false attractor (wrong valley)\\
\textbf{Lock:} EZH2/PRC2 H3K27me3 deposition on
FOXA1, GATA3, ESR1 promoters maintains the wrong
valley state\\
\textbf{FOXA1/EZH2 ratio:} 0.52\\
\textbf{EZH2 elevation:} $+189\%$ above normal
Mature Luminal reference

\medskip
\textbf{Key geometry:}\\
FOXA1 near-absent (silenced by H3K27me3)\\
ESR1 absent (downstream of FOXA1 silencing)\\
KRT5, KRT14, SOX10, FOXC1 expressed
(wrong valley --- basal/myoepithelial identity)\\
EZH2 massively elevated --- this IS the lock\\
The luminal programme is present in the genome.
It is epigenetically silenced. It can be de-silenced.
\vspace{0.5em}
\end{minipage}
}
\end{center}

\subsection{The molecular lock}

EZH2, the catalytic subunit of PRC2, deposits
H3K27me3 (repressive trimethylation) on the promoters
of FOXA1, GATA3, and ESR1. This is not a structural
absence of luminal identity genes. It is a
pharmacologically reversible epigenetic silencing.
Evidence:

\begin{itemize}[noitemsep]
  \item EZH2 $+189\%$ above normal Mature Luminal
        in scRNA-seq data (GSE176078)
  \item FOXA1 near-absent despite intact gene
        (silenced, not deleted)
  \item Toska et al.\ \textit{Nature Medicine} 2017:
        EZH2 inhibition restores FOXA1 expression in
        ER-negative breast cancer
  \item Schade et al.\ \textit{Nature} 2024:
        PRC2/EZH2 inhibition drives TNBC toward
        luminal-like differentiation with GATA3
        induction
  \item The EZH2 paradox: EZH2-high TNBC is
        chemo-sensitive (high proliferation) but
        late-relapsing (H3K27me3 maintains dormant
        residual cell programme). Both arms
        confirmed in GSE25066 ($n=508$; DRFS
        HR$=1.363$, $p=0.0047$).
\end{itemize}

\subsection{The therapeutic unlock sequence}

\begin{center}
\begin{tabular}{p{3.5cm}p{9cm}}
\toprule
\textbf{Step} & \textbf{Action and rationale} \\
\midrule
First step
& \textbf{Tazemetostat} (EZH2 inhibitor, FDA-approved
  2020) to remove H3K27me3 marks from FOXA1, GATA3,
  ESR1 promoters. Required before any downstream agent.
  Without this step, there is no ER to target and no
  luminal identity to engage. \\[0.3em]
Second step\\(conversion therapy)
& \textbf{Fulvestrant} once FOXA1 H-score
  re-expression $\geq50$ is confirmed on Week~4
  biopsy. Engages and degrades the restored ER.
  Full trial proposal: CS-LIT-16 (Zenodo 2026-03-06).
  \\[0.3em]
Maintenance\\(post-chemo window)
& \textbf{Tazemetostat maintenance} after neoadjuvant
  chemotherapy in EZH2-high patients (H-score $\geq150$)
  to eliminate the dormant residual EZH2-high population
  before the 3--5 year late-relapse peak. DRFS at 3 years
  as primary endpoint. Full trial proposal: CS-LIT-17
  (Zenodo 2026-03-06). \\[0.3em]
Patient selection
& EZH2 IHC H-score $\geq150$, FOXA1 IHC H-score
  $<20$ for conversion therapy. EZH2 IHC H-score
  $\geq150$ alone for maintenance. \\
\bottomrule
\end{tabular}
\end{center}

\subsection{The TNBC depth axis within Lock Type 4}

Not all EZH2-high TNBC is at the same depth in the
wrong valley. The framework identifies a continuous
depth axis within TNBC characterised by:

\begin{itemize}[noitemsep]
  \item \textbf{AR as depth axis:}
        AR expression negatively correlates with
        depth score ($r=-0.378$, $p=6.23\times10^{-147}$,
        $n=4{,}312$ cells). Higher AR = shallower
        attractor. LAR (luminal androgen receptor)
        subtype is the shallowest TNBC.
  \item \textbf{TNBC depth score predicts DRFS:}
        HR$=1.509$, $p=0.0001$ in GSE25066 ($n=508$).
        Deeper TNBC = worse distant relapse.
  \item \textbf{Treatment implication:}
        LAR (AR-high) TNBC may respond to androgen
        receptor antagonists (bicalutamide, enzalutamide)
        as a subtype-specific option. The depth score
        as a DRFS predictor is the patient selection
        instrument. (CS-LIT-11, CS-LIT-25)
\end{itemize}

% ══════════════════════════════════════════════════════════════
\section{Lock Type 5: Claudin-low --- The Stem Root Lock}
% ══════════════════════════════════════════════════════════════

\subsection{Attractor geometry}

\begin{center}
\colorbox{mechanismbox}{
\begin{minipage}{0.85\textwidth}
\vspace{0.5em}
\textbf{CLAUDIN-LOW --- TYPE 4 ROOT LOCK}

\medskip
\textbf{Cell of origin:} Mammary stem cell
(not a luminal progenitor)\\
\textbf{Normal terminal attractor:} None --- this
cell never committed to luminal lineage\\
\textbf{Arrest point:} At the progenitor root ---
never descended into any specific lineage valley\\
\textbf{Lock:} Stem cell programme maintenance ---
no luminal, basal, or myoepithelial identity is
present or silenced; it was never activated\\
\textbf{FOXA1/EZH2 ratio:} 0.10 (lowest of all subtypes
in scRNA-seq cancer-cell-specific measurement)\\
\textbf{Luminal programme status:} Absent ---
not silenced, not present

\medskip
\textbf{Key geometry (scRNA-seq, cancer cells only):}\\
FOXA1 at near-zero (lowest of all subtypes)\\
EZH2 low (not the mechanism --- EZH2 marks
proliferating cells; CL is low-proliferation)\\
CDH1, CLDN3, CLDN4, CLDN7 all absent
(defining claudin signature)\\
Stem markers elevated: CD44-high, CD24-low,
ALDH1A3 present\\
Immune compartment dominant: high TIL fraction
\vspace{0.5em}
\end{minipage}
}
\end{center}

\subsection{Why EZH2 inhibition does not apply to CL}

This is the most important distinction between TNBC
(Lock Type 4) and Claudin-low (Lock Type 5). Both
have low FOXA1. But:

\begin{itemize}[noitemsep]
  \item In TNBC: FOXA1 is absent because EZH2 is
        silencing it. The gene is intact. Remove EZH2,
        FOXA1 returns.
  \item In CL: FOXA1 is absent because it was never
        activated. The cell never committed to luminal
        lineage. There is no FOXA1 programme to
        de-silence --- it was never written.
\end{itemize}

Tazemetostat will not restore FOXA1 in CL tumours
because there is nothing to de-repress. The luminal
programme does not exist in these cells in a silenced
form; it does not exist at all. This is the Type 4
vs Type 2 distinction: root lock vs wrong valley.
Treating a CL patient with tazemetostat $\rightarrow$
fulvestrant (the CS-LIT-16 sequence) would be
mechanistically incorrect. The CL patient requires
a completely different therapeutic approach.

\subsection{The molecular lock and its therapeutic
implication}

CL tumours have no epigenetic lock to dissolve and
no kinase amplicon to suppress. The cancer cell
itself is in a stem-like state that is not amenable
to re-differentiation through the locks that work
for Types 1--3. The accessible therapeutic node is
the immune compartment:

\begin{center}
\begin{tabular}{p{3.5cm}p{9cm}}
\toprule
\textbf{Feature} & \textbf{Therapeutic implication} \\
\midrule
High TIL infiltration
& Immune compartment is the primary target.
  Cancer cell re-differentiation is not accessible. \\[0.3em]
FOXP3/CD8A ratio
  HR$=2.212$
& FOXP3/CD8A ratio is the strongest immune OS
  predictor in CL ($p$ confirmed). High FOXP3/CD8A
  (Treg-dominated) predicts worse OS. This is the
  key immune subgroup stratifier. (CS-LIT-20) \\[0.3em]
Memory-low subgroup
& Memory-low CL (Pommier 2020 subgroup 1 proxy)
  is simultaneously deeper ($p=5.52\times10^{-8}$),
  more CT-antigen de-repressed ($p=8.86\times10^{-9}$),
  and more Treg-dominated (FOXP3/CD8A, $p=0.008$).
  Memory-low is the primary anti-TIGIT $+$ anti-PD-1
  beneficiary subgroup. (CS-LIT-19) \\[0.3em]
TIGIT as immune target
& Anti-TIGIT $+$ anti-PD-1 in memory-low CL:
  belrestotug failed in \textit{unselected} TNBC
  (May 2025). Framework predicts it works in
  memory-low CL specifically. SKYLINE trial
  (NCT06175390) has multi-omics biomarker arm.
  The lineage memory score is the proposed selector.
  (CS-LIT-19) \\
\bottomrule
\end{tabular}
\end{center}

\subsection{CL as the deepest subtype: EZH2-free PCA}

Standard PCA including EZH2 as a variable will
place TNBC as the deepest subtype because TNBC
has the highest EZH2 expression. But EZH2 is a
proliferation marker, not a depth marker for CL.
CL is a low-proliferation, stem-like subtype
with low EZH2. When EZH2 is excluded from the
depth PCA (EZH2-free PCA), CL emerges as the
deepest subtype: depth score 6.572 vs TNBC 6.063.
This is the correct measurement because depth
in the Waddington landscape corresponds to
distance from differentiated identity, which
for CL is total (no identity) vs for TNBC is
partial (identity silenced but present).
(CS-LIT-3, CS-LIT-26)

% ══════════════════════════════════════════════════════════════
\section{Lock Type 6: Invasive Lobular Carcinoma ---
The CDH1 Structural Lock}
% ══════════════════════════════════════════════════════════════

\subsection{Attractor geometry}

\begin{center}
\colorbox{mechanismbox}{
\begin{minipage}{0.85\textwidth}
\vspace{0.5em}
\textbf{ILC --- TYPE 3 STRUCTURAL LOCK}

\medskip
\textbf{Cell of origin:} Mature luminal epithelial cell\\
\textbf{Normal terminal attractor:} ER-positive,
organised luminal cell with E-cadherin-mediated
cell-cell adhesion maintaining ductal architecture\\
\textbf{Arrest mechanism:} CDH1 (E-cadherin) loss
removes the structural constraint that maintains
luminal cell polarity and adhesion\\
\textbf{Lock:} Not epigenetic. The lock is the
absence of a structural component (E-cadherin)
that is required for normal luminal cell behaviour
and the maintenance of tissue architecture\\
\textbf{FOXA1/EZH2:} FOXA1 \textit{hyperactivated}
(above LumA levels) --- ILC is the geometric
inverse of TNBC on the identity axis\\
\textbf{Luminal programme:} Fully intact, possibly
hyperactivated

\medskip
\textbf{Key geometry:}\\
CDH1 biallelic loss (mutation $+$ LOH) = defining event\\
FOXA1 amplified or hyperexpressed in a subset\\
ESR1 expressed (ER-positive, endocrine sensitive)\\
VIM absent (not mesenchymal despite single-file growth)\\
The cell is fully luminal. It is growing in single-file
pattern (``Indian file'') because the E-cadherin
scaffold that would organise it into ducts is gone.
This is architectural derangement, not identity loss.
\vspace{0.5em}
\end{minipage}
}
\end{center}

\subsection{Why ILC is the geometric inverse of TNBC}

\begin{center}
\begin{tabular}{p{3.5cm}p{5cm}p{5cm}}
\toprule
\textbf{Feature} & \textbf{TNBC (Lock Type 4)}
& \textbf{ILC (Lock Type 6)} \\
\midrule
Lock mechanism
& Epigenetic (EZH2/H3K27me3)
& Structural (CDH1 loss) \\
FOXA1 status
& Near-absent (silenced)
& Hyperactivated \\
ESR1 status
& Absent (silenced)
& Present (ER$+$) \\
Luminal identity
& Epigenetically locked away
& Structurally unconstrained \\
Lock direction
& Identity programme suppressed
& Identity programme released
  from architectural constraint \\
EZH2 level
& Massively elevated
& Normal or low \\
Therapeutic first step
& Remove the silencing (EZH2i)
& Engage the intact but
  unconstrained ER \\
Resistance pattern
& Endocrine therapy-primary
  resistant
& Aromatase inhibitor-resistant
  (FOXA1 hyperactivation
  overrides AI suppression) \\
\bottomrule
\end{tabular}
\end{center}

\subsection{The molecular lock and therapeutic
implication}

Because ILC has CDH1 loss and FOXA1 hyperactivation:
\begin{itemize}[noitemsep]
  \item Aromatase inhibitors (AIs) are less effective
        in ILC than in IDC: FOXA1 hyperactivation
        amplifies the ER circuit so strongly that
        ligand depletion alone is insufficient to
        suppress output
  \item Fulvestrant (SERD) is superior to AIs in ILC:
        it degrades the ER protein directly rather
        than depleting ligand, bypassing the FOXA1
        amplification
  \item FOXA1 hyperactivation as the stratifier:
        patients with the highest FOXA1 (amplified
        or highest expression) derive the greatest
        differential benefit from fulvestrant over AI
  \item The FOXA1-stratified fulvestrant superiority
        prediction is novel and testable in archived
        tissue from NCT02206984 (FALCON trial). (CS-LIT-18)
\end{itemize}

% ══════════════════════════════════════════════════════════════
\section{The Unified Classification: One Axis, Six Locks}
% ══════════════════════════════════════════════════════════════

\subsection{The complete six lock type table}

\begin{center}
\begin{longtable}{p{1.8cm}p{2cm}p{2.5cm}p{2.5cm}p{4cm}}
\toprule
\textbf{Subtype}
& \textbf{Lock type}
& \textbf{FOXA1/EZH2}
& \textbf{Attractor type}
& \textbf{Therapeutic first step} \\
\midrule
\endhead
Luminal A
& CDK4/6 cycle lock
& 9.38 (highest)
& Type 1
& CDK4/6i $+$ ET \\[0.3em]
Luminal B
& DNMT3A/HDAC2 chromatin lock
& 8.10
& Type 1 (epigenetic brake)
& HDACi (entinostat) then ET \\[0.3em]
HER2-enriched
& ERBB2 kinase override
& 3.34
& Type 1 (kinase arrest)
& Anti-HER2 then ET \\[0.3em]
TNBC/Basal
& EZH2/PRC2 H3K27me3 lock
& 0.52
& Composite 1$\rightarrow$2
& EZH2i (tazemetostat) then fulvestrant
  or maintenance \\[0.3em]
Claudin-low
& Stem root lock
& 0.10 (lowest)
& Type 4
& Anti-TIGIT $+$ anti-PD-1
  (memory-low subgroup) \\[0.3em]
ILC
& CDH1 structural lock
& --- (FOXA1 hyper-activated)
& Type 3
& Fulvestrant (not AI);
  FOXA1 IHC as selector \\
\bottomrule
\caption{The six lock type classification of breast
cancer. FOXA1/EZH2 ratio values from GSE176078
(scRNA-seq per-patient means, 19,542 cancer cells,
26 patients). ILC is not placed on the FOXA1/EZH2
axis because its lock is structural. All therapeutic
first steps are derived from the lock type, not from
receptor status alone.}
\end{longtable}
\end{center}

\subsection{What the classification explains that
receptor status does not}

Current clinical practice classifies breast cancer
by receptor status: ER$+$, PR$+$, HER2$+$, or
triple-negative. This is necessary but insufficient:

\begin{center}
\begin{tabular}{p{4cm}p{4.5cm}p{4.5cm}}
\toprule
\textbf{Clinical question}
& \textbf{Receptor status answer}
& \textbf{Lock type answer} \\
\midrule
Why does entinostat work in some ER$+$ patients
but not others?
& Cannot explain (all ER$+$)
& LumB has chromatin lock; LumA does not.
  TFF1/ESR1 decoupling is the selector. \\[0.3em]
Why does TNBC relapse late despite pCR?
& Cannot explain
& EZH2/PRC2 lock maintains dormant residual cells.
  The EZH2 paradox. Maintenance tazemetostat
  is the intervention. \\[0.3em]
Why does anti-TIGIT fail in unselected TNBC?
& Cannot explain
& TNBC is a wrong-valley type (Lock 4), not a
  root-lock type (Lock 5). Anti-TIGIT is for
  memory-low CL. Wrong subtype. \\[0.3em]
Why does fulvestrant outperform AI in ILC?
& Cannot explain (all ER$+$)
& FOXA1 hyperactivation makes ligand depletion
  insufficient. ER degradation is required.
  Lock type 6 specific. \\[0.3em]
Why is tazemetostat not tested in maintenance?
& No receptor status rationale
& Lock type 4 has a residual EZH2-high dormant
  population after chemo. Maintenance EZH2i
  directly addresses the mechanism. \\
\bottomrule
\end{tabular}
\end{center}

\subsection{What is confirmed vs novel in the classification}

\begin{center}
\begin{tabular}{p{4cm}p{2.8cm}p{6cm}}
\toprule
\textbf{Claim} & \textbf{Status} & \textbf{Evidence} \\
\midrule

CDK4/6 lock in LumA
& CONVERGENT
& Palbociclib, ribociclib, abemaciclib standard of care \\[0.8em]

DNMT3A/HDAC2 coupling in LumB
($r=+0.267$, $p=5.68\times10^{-56}$)
& CONVERGENT-NOVEL
& scRNA-seq confirmed; bulk METABRIC $p=0.0019$ replication \\[0.8em]

TFF1/ESR1 decoupling as HDACi selector
& NOVEL
& CS-LIT-23; NCT07235618 active \\[0.8em]

ERBB2 kinase override as Type~1
& CONVERGENT
& Trastuzumab history \\[0.8em]

EZH2/PRC2 as TNBC lock
& CONVERGENT-NOVEL
& Schade 2024, Toska 2017 \\[0.8em]

Tazemetostat $\rightarrow$ fulvestrant sequence
& NOVEL
& CS-LIT-16 (Zenodo 2026-03-06) \\[0.8em]

Tazemetostat maintenance
& NOVEL
& CS-LIT-17 (Zenodo 2026-03-06) \\[0.8em]

CL as Type~4 (root lock, not Type~2)
& NOVEL
& EZH2-free PCA; scRNA-seq geometry \\[0.8em]

CL memory-low as anti-TIGIT beneficiary
& NOVEL-URGENT
& CS-LIT-19; post-belrestotug field \\[0.8em]

ILC as structural inverse of TNBC
& CONVERGENT-NOVEL
& CDH1 literature; FOXA1 hyperactivation \\[0.8em]

Fulvestrant superiority in ILC by FOXA1 IHC
& NOVEL
& CS-LIT-18; FALCON tissue available \\[0.8em]

FOXA1/EZH2 as single ordering axis
& CONFIRMED\\(7 datasets)
& CS-LIT-1 (Zenodo 2026-03-06) \\[0.4em]

\bottomrule
\end{tabular}
\end{center}
\section{Derivation Provenance and Protocol}
% ══════════════════════════════════════════════════════════════

\subsection{The before-document protocol}

All six lock type assignments were made using the
before-document protocol: each assignment was locked
in a before-document prior to any script running on
the relevant dataset. The document chain for each
subtype:

\begin{center}
\begin{tabular}{p{2.5cm}p{3cm}p{7cm}}
\toprule
\textbf{Subtype} & \textbf{Before-document}
& \textbf{Repository path} \\
\midrule
LumA   & BRCA-S1a & \texttt{.../LumA/} \\
LumB   & BRCA-S5a & \texttt{.../LumB/} \\
HER2   & BRCA-S3a & \texttt{.../HER2\_ENRICHED/} \\
TNBC   & BRCA-S4a & \texttt{.../TNBC/predictions.md} \\
CL     & BRCA-S7a & \texttt{.../Claudin-low/} \\
ILC    & BRCA-S6a & \texttt{.../ILC/} \\
Cross  & BRCA-S8h & \texttt{.../BRCA\_Cross\_Subtype\_}\\
       &          & \texttt{Literature\_Check.md} \\
\bottomrule
\end{tabular}
\end{center}

All documents are version-controlled with cryptographic
timestamps at:
\href{https://github.com/Eric-Robert-Lawson/attractor-oncology}
{\texttt{github.com/Eric-Robert-Lawson/attractor-oncology}}

\subsection{Single-cell data: GSE176078}

All ratio values and depth scores are derived from
GSE176078 (Wu et al.\ \textit{Nature Genetics} 2021):
\begin{itemize}[noitemsep]
  \item 100,064 total cells; 19,542 cancer cells
        used for ratio and depth derivation
  \item 26 patients; 4 PAM50 subtypes in scRNA-seq
        ($+$ claudin-low and ILC from TCGA bulk)
  \item Per-patient mean FOXA1 / mean EZH2 computed
        across all cancer cells from each patient
  \item Scripts: \texttt{ratio1.py--ratio4.py} in
        \texttt{.../FOXA1\_EZH2\_RATIO/}
\end{itemize}

% ══════════════════════════════════════════════════════════════
\section{The Unifying Biological Principle}
% ══════════════════════════════════════════════════════════════

\begin{center}
\colorbox{summarybox}{
\begin{minipage}{0.88\textwidth}
\vspace{0.6em}
\textbf{THE UNIFYING PRINCIPLE}

\medskip
Breast cancer subtypes are not a random collection
of molecular accidents. They are a structured set
of positions in a single developmental landscape,
each maintained by a specific molecular lock.

\medskip
The FOXA1/EZH2 ratio encodes position on the luminal
identity axis. High ratio: the cell is fully luminal
and its lock is in the cycle or chromatin machinery.
Low ratio: the cell has lost luminal identity through
epigenetic silencing (TNBC) or never had it (CL).
ILC is the structural exception: FOXA1 is
hyperactivated but E-cadherin is gone.

\medskip
Each lock type has exactly one correct first step:
\begin{itemize}[noitemsep,leftmargin=*]
  \item LumA: CDK4/6 inhibitor (cycle lock)
  \item LumB: HDACi (chromatin brake)
  \item HER2: anti-HER2 (kinase override)
  \item TNBC: EZH2 inhibitor (epigenetic silencing)
  \item CL: anti-TIGIT $+$ anti-PD-1 (immune only)
  \item ILC: fulvestrant, not AI (structural hyper-FOXA1)
\end{itemize}

\medskip
Applying the wrong first step is not merely
suboptimal. In some cases it is actively incorrect:
\begin{itemize}[noitemsep,leftmargin=*]
  \item Giving endocrine therapy to TNBC without
        EZH2i first: no ER to engage (silenced)
  \item Giving tazemetostat to CL: no FOXA1 programme
        to de-silence (never written)
  \item Giving AI to FOXA1-hyperactivated ILC:
        ligand depletion is insufficient when the
        ER amplifier is structurally unconstrained
  \item Giving anti-TIGIT to unselected TNBC:
        wrong subtype (TNBC is Type 2, CL is Type 4)
\end{itemize}

\medskip
The classification makes the sequencing logic explicit.
Each treatment decision follows from the lock type,
not from the receptor status alone.
\vspace{0.5em}
\end{minipage}
}
\end{center}

% ══════════════════════════════════════════════════════════════
\section{Locked Statement}
% ══════════════════════════════════════════════════════════════

\begin{center}
\colorbox{novelbox}{
\begin{minipage}{0.88\textwidth}
\vspace{0.6em}
\textbf{LOCKED AS OF 2026-03-06}

\medskip
The six lock type classification of breast cancer is
derived from Waddington landscape attractor geometry
applied to 19,542 single cancer cells before literature
review. All six lock types have published biological
support. Three generate novel trial proposals
(CS-LIT-16, CS-LIT-17, CS-LIT-18, CS-LIT-19,
CS-LIT-21, CS-LIT-23) with no currently registered
trials.

\medskip
The classification is ordered by the FOXA1/EZH2 ratio
(validated in $\sim$7,500 patients across seven
independent datasets, companion Zenodo deposit
2026-03-06). It provides a mechanistically unified
explanation for why these six subtypes exist, why they
require different therapeutic first steps, and why
applying those steps out of order fails.

\medskip
No competing attractor geometry classification of
breast cancer subtypes currently exists in the
published literature.

\medskip
\hfill --- Eric Robert Lawson, March 6, 2026
\vspace{0.4em}
\end{minipage}
}
\end{center}

% ══════════════════════════════════════════════════════════════
\section*{Repository and Reproducibility}
\addcontentsline{toc}{section}{Repository and
Reproducibility}
% ══════════════════════════════════════════════════════════════

\begin{center}
\href{https://github.com/Eric-Robert-Lawson/attractor-oncology}
{\texttt{https://github.com/Eric-Robert-Lawson/attractor-oncology}}\\[0.3em]
\small Repository ID: 1172048282 \textbar{}
MIT License \textbar{}
Python \textbar{}
Public
\end{center}

\begin{center}
\begin{tabular}{ll}
\toprule
\textbf{Content} & \textbf{Path} \\
\midrule
Attractor geometry axioms
& \texttt{ATTRACTOR\_GEOMETRY\_AXIOMS.md} \\
BRCA subtype orientation
& \texttt{Cancer\_Research/BRCA/DEEP\_DIVE/}
  \texttt{BRCA\_Subtypes.md} \\
Cross-subtype literature check
& \texttt{.../BRCA\_Cross\_Subtype\_}
  \texttt{Literature\_Check.md} \\
LumA deep dive & \texttt{.../LumA/} \\
LumB deep dive & \texttt{.../LumB/} \\
HER2 deep dive & \texttt{.../HER2\_ENRICHED/} \\
TNBC deep dive & \texttt{.../TNBC/} \\
Claudin-low deep dive & \texttt{.../Claudin-low/} \\
ILC deep dive & \texttt{.../ILC/} \\
FOXA1/EZH2 ratio scripts
& \texttt{.../FOXA1\_EZH2\_RATIO/ratio1--4.py} \\
\bottomrule
\end{tabular}
\end{center}

% ══════════════════════════════════════════════════════════════
\section*{Author Statement}
\addcontentsline{toc}{section}{Author Statement}
% ══════════════════════════��═══════════════════════════════════

This work was conducted independently by Eric Robert
Lawson under the OrganismCore research programme. No
institutional funding was received. No conflicts of
interest exist. All data used is publicly available.
All code is open source under the MIT licence.

% ══════════════════════════════════════════════════════════════
\section*{Contact}
\addcontentsline{toc}{section}{Contact}
% ══════════════════════════════════════════════════════════════

\begin{tabular}{ll}
Author        & Eric Robert Lawson \\
Organisation  & OrganismCore \\
Email         &
  \href{mailto:OrganismCore@proton.me}
       {OrganismCore@proton.me} \\
ORCID         &
  \href{https://orcid.org/0009-0002-0414-6544}
       {https://orcid.org/0009-0002-0414-6544} \\
Repository    &
  \href{https://github.com/Eric-Robert-Lawson/attractor-oncology}
       {github.com/Eric-Robert-Lawson/attractor-oncology} \\
Homepage      &
  \href{https://github.com/Eric-Robert-Lawson/OrganismCore}
       {github.com/Eric-Robert-Lawson/OrganismCore} \\
Date          & 2026-03-06 \\
\end{tabular}

\vspace{2em}

\begin{center}
\itshape
``Six subtypes. Six locks. Six keys.\\
The same axis orders them all.\\
The FOXA1/EZH2 ratio is the ruler.\\
The lock type is the prescription.''

\medskip
\upshape
--- Eric Robert Lawson, March 6, 2026
\end{center}

% ══════════════════════════════════════════════════════════════
\end{document}
